\documentclass[border=10pt,tikz]{standalone}
\usepackage{fontspec}
\setmainfont{Apple SD Gothic Neo}
\setsansfont{Apple SD Gothic Neo}
\usepackage{tikz}
\usetikzlibrary{positioning,arrows.meta,shapes,fit,backgrounds,calc,decorations.pathreplacing}
\begin{document}
% _design.tex — shared TikZ design system for all blog diagrams
% Inputs nothing on its own — meant to be \input{} from each diagram .tex
% AFTER \begin{document}, BEFORE \begin{tikzpicture}.
%
% Provides a unified visual language across GoF / DSA / EMC++ / Beautiful series:
%   - Colors: muted, accessible, color-coded by role
%   - Node styles: consistent sizing, rounded, subtle borders
%   - Edge styles: UML-style inheritance, aggregation, plus dashed alternates
%
% Usage:
%   \documentclass[border=10pt,tikz]{standalone}
%   \usepackage{tikz}
%   \usetikzlibrary{positioning,arrows.meta,shapes,fit,backgrounds}
%   \begin{document}
%   \input{../_design.tex}     % adjust relative path
%   \begin{tikzpicture}[blog]  % apply 'blog' style
%   ...

\definecolor{absfill}{RGB}{220,232,247}   % cool blue
\definecolor{absbord}{RGB}{72,118,182}
\definecolor{confill}{RGB}{249,222,222}   % warm red
\definecolor{conbord}{RGB}{197,93,93}
\definecolor{impfill}{RGB}{223,238,222}   % soft green
\definecolor{impbord}{RGB}{99,154,93}
\definecolor{clifill}{RGB}{251,239,205}   % pale yellow
\definecolor{clibord}{RGB}{200,167,72}
\definecolor{hifill} {RGB}{253,224,200}   % accent orange
\definecolor{hibord} {RGB}{210,135,68}
\definecolor{nodefill}{RGB}{250,250,250}
\definecolor{nodebord}{RGB}{120,120,120}
\definecolor{edgec}  {RGB}{80,80,80}
\definecolor{edgesc} {RGB}{145,145,145}
\definecolor{groupbord}{RGB}{200,200,200}

\tikzset{
  blog/.style={
    font=\sffamily\small,
    every node/.style={font=\sffamily\small}
  },
  % --- node roles ---
  cls/.style={
    rectangle, rounded corners=2pt,
    draw=nodebord, line width=0.4pt,
    minimum height=8mm, minimum width=28mm,
    inner sep=3pt, fill=nodefill, align=center
  },
  abs/.style={cls, draw=absbord, fill=absfill, font=\sffamily\itshape\small},
  con/.style={cls, draw=conbord, fill=confill},
  imp/.style={cls, draw=impbord, fill=impfill},
  cli/.style={cls, draw=clibord, fill=clifill},
  hi/.style={cls, draw=hibord, fill=hifill, font=\sffamily\bfseries\small},
  % --- circle (for trees / graph nodes) ---
  vx/.style={
    circle, draw=absbord, line width=0.4pt,
    minimum size=8mm, fill=absfill, inner sep=0pt
  },
  vxc/.style={vx, draw=conbord, fill=confill},
  vxh/.style={vx, draw=hibord, fill=hifill, font=\sffamily\bfseries},
  % --- decision diamond ---
  q/.style={
    diamond, draw=clibord, line width=0.4pt, fill=clifill,
    inner sep=2pt, aspect=2.4, align=center
  },
  % --- edges ---
  inh/.style={                                  % UML inheritance
    -{Triangle[length=3mm,width=3mm,open]},
    draw=edgec, line width=0.5pt
  },
  agg/.style={                                  % UML aggregation
    {Diamond[length=2.5mm,width=2.5mm,open]}-{Stealth[length=2.2mm]},
    draw=edgec, line width=0.45pt
  },
  uses/.style={                                 % directed reference
    -{Stealth[length=2.2mm]}, draw=edgec, line width=0.5pt
  },
  arr/.style={                                  % plain arrow (graphs/trees)
    -{Stealth[length=2mm]}, draw=edgec, line width=0.45pt
  },
  edge/.style={                                 % undirected (trees)
    draw=edgec, line width=0.45pt
  },
  alt/.style={                                  % dashed: similar/alt
    -{Stealth[length=2mm]}, dashed, draw=edgesc, line width=0.45pt
  },
  ret/.style={                                  % dashed creates/returns
    -{Stealth[length=2mm]}, dashed, draw=edgesc, line width=0.45pt
  },
  thr/.style={                                  % threading link (red dashed)
    ->, dashed, draw=conbord, line width=0.5pt
  },
  % --- group / region ---
  group/.style={
    rectangle, rounded corners=4pt,
    draw=groupbord, line width=0.4pt,
    inner sep=10pt, fill=black!2
  },
  glabel/.style={font=\sffamily\bfseries\small, anchor=north west, inner sep=2pt}
}

% Korean font convenience macro — included by every Hangul diagram
\providecommand{\KOR}{}

\begin{tikzpicture}[blog, flowb/.style={cls, text width=2.9cm, minimum height=12mm}]
\node[font=\sffamily\bfseries\large] at (5.4,2.5) {필드 불량에서 근본 원인까지};
\node[flowb, neut] (a0) at (0.0,0) {\textbf{필드 불량 보고}};
\node[flowb, con] (a1) at (3.6,0) {\textbf{로그·ramdump 수집}};
\draw[arr] (a0.east) -- (a1.west);
\node[flowb, prp] (a2) at (7.2,0) {\textbf{통계로 분류}};
\draw[arr] (a1.east) -- (a2.west);
\node[flowb, hi] (a3) at (10.8,0) {\textbf{불량 샘플 회수 (RMA)}};
\draw[arr] (a2.east) -- (a3.west);
\node[flowb, prp] (b0) at (10.8,-2.4) {\textbf{재현 시험}\\[2pt]\footnotesize 시스템·테스터에서};
\node[flowb, imp] (b1) at (7.2,-2.4) {\textbf{비파괴 분석}\\[2pt]\footnotesize X-ray·초음파 현미경};
\draw[arr] (b0.west) -- (b1.east);
\node[flowb, warn] (b2) at (3.6,-2.4) {\textbf{파괴 분석}\\[2pt]\footnotesize 디캡·단면 절단};
\draw[arr] (b1.west) -- (b2.east);
\node[flowb, ok] (b3) at (0.0,-2.4) {\textbf{근본 원인과 조치}\\[2pt]\footnotesize 설계·공정·SW};
\draw[arr] (b2.west) -- (b3.east);
\draw[arr] (a3.south) -- (b0.north);
\draw[decorate, decoration={brace, amplitude=5pt, mirror}, draw=mutec] (1.9,-3.3) -- (8.9,-3.3);
\node[caption, anchor=north] at (5.4,-3.55) {FA(실패 분석) 팀의 영역};
\draw[decorate, decoration={brace, amplitude=5pt}, draw=mutec] (-1.5,0.85) -- (8.7,0.85);
\node[caption, anchor=south] at (3.6,1.0) {SW·시스템 팀이 데이터를 모으는 영역};
\end{tikzpicture}
\end{document}
